\documentclass[11pt,a4paper,sans]{moderncv}
\moderncvstyle{classic}
\moderncvcolor{blue}

\usepackage[utf8]{inputenc}
\usepackage[brazilian]{babel}
\usepackage[scale=0.82]{geometry}
\usepackage{hyperref}

\name{Marcelo}{Vironda Rozanti}
\title{Engenheiro de Software Backend Sênior}
\address{Rua Piauí 163}{São Paulo - SP, Brasil}{}
\email{mvrozanti@hotmail.com}
\homepage{mvrozanti.github.io}
\social[linkedin]{mvrozanti}
\social[github]{mvrozanti}

\begin{document}
\makecvtitle

\section{Resumo}
Engenheiro backend JVM com nove anos em bancos e fintechs brasileiras. Confortável em assumir serviços críticos e regulados do desenho à operação. Foco recente: sistemas de decisão de crédito, pagamentos instantâneos, automação fiscal.

\section{Experiência}

\cventry{04/2025 -- atual}{Engenheiro Sênior}{C6 Bank}{São Paulo}{}{%
Responsável pelo desenho, implementação e operação de serviços críticos de alta disponibilidade para concessão de crédito. Codebase Kotlin com Domain-Driven Design; observabilidade e padrões de recuperação como cidadãos de primeira classe.
\begin{itemize}
\item Kotlin, Spring, MongoDB, PostgreSQL, Docker
\item Harness em CI/CD; Splunk e Grafana para observabilidade
\end{itemize}}

\cventry{03/2024 -- 04/2025}{Engenheiro de Software}{Fiserv}{São Paulo}{}{%
Serviços da plataforma POS para pagamentos presenciais. Responsável pela postura de prevenção e recuperação de desastres; liderou a adoção de DevSecOps.
\begin{itemize}
\item Spring, Oracle SQL, Docker, Redis
\item Splunk + Dynatrace; Fortify, Sonar, CyberArk em CI
\end{itemize}}

\cventry{09/2021 -- 08/2022}{Desenvolvedor Backend}{Conta Azul}{São Paulo}{}{%
Construiu automação de obrigações fiscais brasileiras (SPED, eSocial, EFD-Reinf) para clientes contábeis SMB. Primitivos AWS-nativos, IaC via Terraform.}

\cventry{12/2020 -- 09/2021}{Desenvolvedor Backend}{Grão}{São Paulo}{}{%
Implementou a integração com o Pix em Kotlin, do kickoff ao lançamento. Ambientes containerizados; modelagem focada em performance; APIs REST escaláveis.}

\cventry{01/2020 -- 12/2020}{Desenvolvedor}{Raia Drogasil}{São Paulo}{}{%
Adequação à LGPD do software de Operações de Venda. Integração de leitores biométricos e impressoras térmicas; modelagem diária em Oracle.}

\cventry{03/2016 -- 01/2020}{Desenvolvedor}{Sys4Bank}{São Paulo}{}{%
Serviços de back-office bancário sobre Spring Boot. Introduziu containerização e Travis CI; modelagem diária em MSSQL.}

\section{Formação}
\cventry{2020}{Bacharel em Ciência da Computação}{Universidade Presbiteriana Mackenzie}{São Paulo}{}{%
Monografia: ``Três experimentos com escalonamento assíncrono de atualização por prioridade de vizinhança em autômatos celulares elementares''.}

\section{Open source}
\cvitem{\href{https://github.com/mvrozanti/RAT-via-Telegram}{RAT-via-Telegram}}{Windows RAT via Telegram (600+ estrelas)}
\cvitem{\href{https://github.com/mvrozanti/dte}{dte}}{Linguagem de processamento de data \& hora (300+ estrelas)}
\cvitem{\href{https://github.com/mvrozanti/dotty}{dotty}}{Versionamento de dotfiles}

\section{Habilidades}
\cvitem{Linguagens}{Java, Kotlin, Python, C/C++, Shell}
\cvitem{Frameworks}{Spring, Spring Boot, Javalin, Koin}
\cvitem{Dados}{PostgreSQL, MongoDB, MySQL, Oracle, MSSQL, Redis}
\cvitem{Plataforma}{Docker, AWS, Terraform, GitLab CI, Harness, Travis CI}
\cvitem{Idiomas}{Português (nativo), Inglês (fluente), Espanhol (intermediário)}

\end{document}
